\documentclass[8pt]{beamer}

\usetheme{Copenhagen}
\usecolortheme{seahorse}
\setbeamertemplate{headline}{}

\usepackage{amsmath, amsthm, amssymb}
\usepackage{mathpartir}
\usepackage{extarrows}
\usepackage{stmaryrd}
\usepackage{xcolor}
\usepackage{naive-ebnf}
\usepackage{bm}
\usepackage{url}
\usepackage{hyperref}

\definecolor{sblue}{RGB}{41, 128, 185}

\usepackage{pl}
\usepackage{iris}

\title{Modeling Pseudo-randomness in Guarded Intearction Trees}
\author{Charles Peng}
\date{\today}

\AtBeginSection[]
{
  \begin{frame}<beamer>
    \frametitle{Plan}
    \tableofcontents[currentsection]
  \end{frame}
}

\newcommand{\prng}{\mathrm{prng}}
\newcommand{\subeff}{\rightarrowtail}

\newcommand{\Error}{\mathrm{Error}}
\newcommand{\Ins}{\mathsf{Ins}}
\newcommand{\Outs}{\mathsf{Outs}}

\newcommand{\new}{\mathrm{new}}
\newcommand{\del}{\mathrm{del}}
\newcommand{\seed}{\mathrm{seed}}
\newcommand{\rand}{\mathrm{rand}}

\newcommand{\efinput}{\mathrm{input}}
\newcommand{\efoutput}{\mathrm{output}}

\newcommand{\efalloc}{\mathrm{alloc}}
\newcommand{\efdealloc}{\mathrm{dealloc}}
\newcommand{\efread}{\mathrm{read}}
\newcommand{\efwrite}{\mathrm{write}}

\newcommand{\redu}{\mapsto}
\newcommand{\redus}{\redu^\ast}
\newcommand{\Redu}{\rightsquigarrow}
\newcommand{\Redus}{\Redu^\ast}
\newcommand{\easgn}{\;\syn{\mathrel{:=}}\;}

\newcommand{\Set}{\mathbf{Set}}
\newcommand{\OFE}{\mathbf{OFE}}
\newcommand{\COFE}{\mathbf{COFE}}

\newcommand{\inl}{\mathrm{inl}}
\newcommand{\inr}{\mathrm{inr}}

\newcommand{\hasstate}{\textsf{has\_state}}

\newcommand{\IT}{\mathrm{IT}}
\newcommand{\gRet}{\mathrm{Ret}}
\newcommand{\gErr}{\mathrm{Err}}
\newcommand{\gTau}{\mathrm{Tau}}
\newcommand{\gVis}{\mathrm{Vis}}
\newcommand{\gFun}{\mathrm{Fun}}
\newcommand{\Tick}{\mathrm{Tick}}
\DeclareMathOperator*{\Next}{\mathrm{next}}
\newcommand{\getret}{\mathrm{get\_ret}}
\newcommand{\getfun}{\mathrm{get\_fun}}
\newcommand{\getval}{\mathrm{get\_val}}
\newcommand{\cbnap}{\bullet_{\text{cbn}}}
\newcommand{\cbvap}{\bullet_{\text{cbv}}}
\newcommand{\gnatop}[1]{\mathrm{NatOp}_{#1}}
\newcommand{\gif}{\mathrm{If}}
\newcommand{\gwhile}{\mathrm{While}}

\newcommand{\out}{\mathrm{read}}
\newcommand{\upd}{\mathrm{update}}
\newcommand{\mix}{\mathrm{mix}}
\newcommand{\knownprng}{\textsf{known\_prng}}
\newcommand{\hasprngs}{\textsf{has\_prngs}}
\newcommand{\prngCtx}{\knowInv{\iota p}{\exists \sigma.\; \hasstate(\sigma) * \hasprngs(\sigma)}}
\newcommand{\heapCtx}{\knowInv{\iota h}{\exists \sigma_h.\; \hasstate(\sigma_h)}}

\newcommand{\Excl}{\operatorname{Excl}}
\newcommand{\AuthFull}{\operatorname{\bullet}\;}
\newcommand{\AuthFrag}{\operatorname{\circ}\;}
\newcommand{\EmptyK}{\mathrm{EmptyK}}
\newcommand{\recval}[3]{\syn{rec}\;#1(#2)=#3}

\newcommand{\demph}[1]{\textcolor{gray}{#1}}

\begin{document}
  \maketitle{}

  \begin{frame}{Overview}

    \begin{block}{Efforts and Contributions}
      \begin{enumerate}
        \item (Background) \demph{Iris in Rocq, SKLR in Iris, Model of SGDT\&Iris, GItrees, etc.}
        \item PRNG effect theory \demph{signature, reifier, and program logic.}
        \item GItree-based semantics \demph{computational adequacy for \(\lambda_\prng\)}
        \item Formalization in Rocq as an extension to the GItree development
        \item Discussion of \demph{limitations and possible generalizations}
      \end{enumerate}
    \end{block}

    \pause

    \begin{block}{Style guide}
      \begin{itemize}
        \item mostly high-level and informal
        \item focus on what I have actually done
        \item non-standard and yet non-original
      \end{itemize}
    \end{block}
  \end{frame}

  \section{Iris --- a logic of recursion}

  \begin{frame}{Fixed point in Iris}
    From \emph{topos of tree} (SGDT) to the category of \textbf{OFE/COFE}.
    \begin{block}{induction on recursive structures --- iProp level iso-recursion}
      Guarded recursion: Banach's fixed point theorem. (Contractive function)
      \begin{mathpar}
        \inferrule[t-fix]{\typed{\Gamma,x:\tau}{t}{\tau} \\ \text{later-guarded}}{\typed{\Gamma}{\mu x . t}{\tau}}
        \and \inferrule[eq-fix]{ }{\typed{\Gamma}{\mu x .\; t = \subst{t}{x}{(\mu x .\; t)}}{\tau}}
        \and \inferrule[l\"ob]{ \later P \vdash P }{\vdash P}
      \end{mathpar}
      Examples: recursive representation predicates and weakest preconditions.
    \end{block}
    \begin{block}{constructing recursive structures --- RDE solver}
      America and Rutten's theorem. (Locally contractive bi-functor on \textbf{COFE}. Non-empty.)
      \[
        G : \COFE^{op}\times \COFE \to \COFE
        \qquad
        G(T, T) \cong T
      \]
      Examples: Iris model and guarded interaction trees
    \end{block}
  \end{frame}

  \begin{frame}{SKLR in Iris, abstractly}
    \begin{itemize}
      \item Step-indexed interpretations of iProp. iProp monotone w.r.t. step-indices.
      \item Kripke frame: Worlds \( \ownM{\cdot}\). Accessibility \(\subseteq = \ownM{a} \vdash \pvs \ownM{b}\) when \(a\Redu b\).\\
        iProp monotone w.r.t. resources.
      \item MISC: iProp impredicativity. Guarded recursion. WP closure under expansion.
      \item Persistent value relation: contraction* and weakening in object language.
    \end{itemize}
    \color{gray}
    \begin{align*}
      \semEtype{\Tenv}{\typ}{\semenv}(\expr) \eqdef{} 
      & \wpre{e}{\semVtype{\Tenv}{\typ}{\semenv}}\\
      \semVtype{\Tenv}{\tvar}{\semenv}(\val) \eqdef{}
      & \semenv(\tvar)(\val)\\
      \semVtype{\Tenv}{\Tnat}{\semenv}(\val) \eqdef{}
      & \exists n \in \nat.\; \val = n\\
      \semVtype{\Tenv}{\typ_1 \times \typ_2}{\semenv}(\val) \eqdef{}
      & \exists \val_1, \val_2.\; \val = (\val_1, \val_2) \sep
      \semVtype{\Tenv}{\typ_1}{\semenv}(\val_1) \sep \semVtype{\Tenv}{\typ_2}{\semenv}(\val_2)\\
      \semVtype{\Tenv}{\typ \to \typ'}{\semenv}(\val) \eqdef{}
      & \persistently \left(\forall \valB.\; \semVtype{\Tenv}{\typ}{\semenv}(\valB) \wand 
      \semEtype{\Tenv}{\typ'}{\semenv}(\val~\valB) \right) \\
      \semVtype{\Tenv}{\forall \tvar.\; \typ}{\semenv}(\val) \eqdef{}
      & \persistently\left(\forall (\predB : \iVal \to \iProp_\persistently).\; \semEtype{\Tenv, \tvar}{\typ}{\semenv, \tvar \mapsto \predB}(\val~\langle\rangle)\right)\\
      \semVtype{\Tenv}{\Tmu \tvar.\; \typ}{\semenv}(\val) \eqdef{}
      & \mu (\predB : \iVal \to \iProp_\persistently).\; \exists \valB.\; (\val = \textlang{fold}\ \valB) \sep
      \later \semVtype{\Tenv}{\typ}{\semenv, \tvar \mapsto \predB}(\valB) \\
      \semVtype{\Tenv}{\Tref(\typ)}{\semenv}(\val) \eqdef{}
      & \exists \loc.\; \val = \loc \sep
      \knowInv{\namesp.\loc}{\exists \valB.\; \loc \mapsto \valB \sep
      \semVtype{\Tenv}{\typ}{\semenv}(\valB)}
    \end{align*}
  \end{frame}

  \begin{frame}{Interlude: modularity of proofs parameterized by CMRAs}
    \begin{block}{data types \`a la carte}
      \begin{mathpar}
        \inferrule{x \in A}{\exists B.\; A \cong B + x}
        \and \inferrule{ }{x \in x\col A}
        \and \inferrule{x \in A}{x \in A+B}
        \and \inferrule{x \in B}{x \in A+B}
      \end{mathpar}
      Universally quantify over an open union \(U\). Contextual/Implicit variable \(x \in A\)\\
      In the end, instantiate \(U\) with a big enough concrete sum.
    \end{block}

    \begin{exampleblock}{}
      \begin{itemize}
        \item CMRA: \texttt{gFunctors} and \texttt{subG}
        \item Reifier: \texttt{gReifiers} and \texttt{subReifier}
        \item Base value: \texttt{sumO} and \texttt{subOfe}
      \end{itemize}
    \end{exampleblock}
  \end{frame}

  \section{GItree --- a domain for untyped effectful programs}

  \begin{frame}{effects signatures}
    \begin{exampleblock}{we need your signatures}
      To specify an effect call: \(\gVis_i(x, k)\)
      \begin{itemize}
        \item operator \(i\)
        \item input data \(x\)
        \item callback/continuation \(k\)
      \end{itemize}
      Signature: defining well-formed calls \((E : \Set,\Ins, \Outs : E \to \COFE \to \COFE)\)
    \end{exampleblock}
    \pause
    First-order effects: single tape I/O \(E = \{\efinput, \efoutput\}\)
    \[
      \Ins_i(X) = \begin{cases}
        \mathbf{1} & i = \efinput\\
        \mathbb{N} & i = \efoutput\\
      \end{cases}
      \quad
      \Outs_i(X) = \begin{cases}
        \mathbb{N} & i = input\\
        \mathbf{1} & i = output\\
      \end{cases}
    \]
    Higher-order effects: untyped reference cells \(E = \{\efalloc, \efdealloc, \efread, \efwrite\}\)
    \[
      \Ins_i(X) = \begin{cases}
        X & i = \efalloc\\
        \iLoc & i = \efdealloc\\
        \iLoc & i = \efread\\
        \iLoc \times X & i = \efwrite\\
      \end{cases}
      \quad
      \Outs_i(X) = \begin{cases}
        \iLoc & i = \efalloc\\
        \mathbf{1} & i = \efdealloc\\
        X & i = \efread\\
        \mathbf{1} & i = \efwrite\\
      \end{cases}
    \]
  \end{frame}

  \begin{frame}[allowframebreaks]{a domain of higher-order effectful programs}
    Fixing the set of base values \(A\) and a effect signature:
    \begin{align*}
      D & \cong \Error + A + \latert D + \latert [D \to D] + \sum_{i\in E} \left( \Ins_i(\latert D) \times [\Outs_i(\latert D) \to \latert D] \right) \cong F(D,D) \\
      F(X,Y) & = \Error + A + \latert Y + \latert [X \to Y] + \sum_{i\in E} \left( \Ins_i(\latert Y) \times [\Outs_i(\latert X) \to \latert Y] \right)
    \end{align*}

    Fold/Unfold + Guarded Rec: intro\&elim\&comp untyped effectful calculus
    \begin{block}{Constructors}
      \begin{itemize}
        \item Error \(\gErr(e : \Error)\)
        \item Base value \(\gRet(a : A)\)
        \item Function value \(\gFun(f : \latert (\IT \to \IT))\)
        \item Delay \(\gTau(t : \latert \IT)\). Normal form \(\Tick(\alpha : \IT) \eqdef \gTau(\Next(\alpha))\)
        \item Visible event \(\gVis_{i \in E}(x: \Ins_i(\latert \IT),\,k : \Outs_i(\latert \IT) \to \latert \IT)\)
      \end{itemize}
    \end{block}

    \framebreak

    \begin{block}{Homomorphisms}
      A homomorphism \(f : \IT \to \IT\) cares only about GItree values.
      \[
        f(\gErr(e)) \equiv \gErr(e)
        \quad
        f(\Tick(\alpha)) \equiv \Tick(f(\alpha))
        \quad
        f(\gVis_i(x,k)) \equiv \gVis_i (x, \latert f \circ k)
      \]
      \begin{itemize}
        \item extract base value \(\getret(f,-)\): \(\getret(f,\gRet(a)) \equiv f\,a\)
        \item extract function \(\getfun(g,-)\): \(\getfun(g,\gFun(h)) \equiv g\,h\)
        \item CBV fun app \(- \cbvap \beta\): \(\gFun(\Next f) \cbvap \beta \equiv f(\beta)\)
        \item Control flow: let-in, sequencing, if-conditional, while-loops, etc.
      \end{itemize}
    \end{block}
  \end{frame}

  \begin{frame}{GItree operational semantics}
    Equational theory --- pure computation. Reifiers --- effectful reduction.\\
    \begin{block}{Effect reification}
      Effect state type and state transition function
      \[
        \State : \COFE \to \COFE
        \quad
        r : \prod_{i\in E} \left( \Ins_i(X) \times \State(X) \to \option{(\Outs_i(X) \times \State(X))} \right)
      \]
      (One possible) I/O effect implementation \(\State(X) = \mathbb{N}^\ast \times \mathbb{N}^\ast\)
      \[
        \quad
        r_{0}((),(l_0,l_1)) = \begin{cases}
          \None & l_0 = [] \\
          \Some(x, (l_0',l_1)) & l_0 = x \col l_0'
        \end{cases}
        \quad
        r_{1}(x,(l_0,l_1)) = \Some((), (l_0, x \col l_1))
      \]
    \end{block}
    \begin{block}{Running a GItree program}
      Configuration: \(\IT \times \State(\IT)\)
      \begin{enumerate}
        \item unfolding guarded fixed points and applying functions \((\Tick(\alpha),\sigma) \Redu (\alpha,\sigma)\).
        \item executing effects \((\gVis_i(x, k),\sigma) \Redu (\beta,\sigma)\) when \(r_i(x) = \Some(y,\sigma')\) and \(k\ y = \Next \beta\).
      \end{enumerate}
      GItree homomorphisms are evaluation contexts.
    \end{block}
  \end{frame}

  \begin{frame}{GItree program logic}
    The guarded fixed point definition of weakest precondition. \(\wpre{\alpha}{\Phi}\)
    \begin{itemize}
      \item Either \(\alpha_v\) and \(\Phi(\alpha_v)\).
      \item Or \((\alpha,\sigma) \Redu (\alpha',\sigma')\) and \(\wpre{\alpha'}{\Phi}\), for every \(\sigma,\sigma'\).
        % Use ghost states and frame-preserving updates instead of actual reduction.
    \end{itemize}
    The familiar proof rules
    \begin{mathpar}
      \inferrule[wp-val]{\Phi(\alpha_v)}{\wpre{\alpha_v}{\Phi}}
      \and
      \inferrule[wp-tick]{\later\wpre{\alpha}{\Phi}}{\wpre{\Tick(\alpha)}{\Phi}}
      \and
      \inferrule[wp-bind]{\text{IT-Hom}(f) \\ \wpre{\alpha}{\alpha_v . \wpre{f(\alpha_v)}{\Phi}}}{\wpre{f(\alpha)}{\Phi}}
      \and
      \inferrule[wp-upd]{\pvs \wpre{\alpha}{\alpha_v . \pvs \Phi(\alpha_v)}}{\wpre{\alpha}{\Phi}}
      \and
      \inferrule[wp-mono]{\wpre{\alpha}{\Psi} \\ \forall \alpha_v . \Psi(\alpha_v)\wand\Phi(\alpha_v)}{\wpre{\alpha}{\Phi}}
      \and
      \inferrule[wp-reify]
      {\hasstate(\sigma)\\
        r_i(x,\sigma) = \Some(\beta,\sigma')\\
      \later \left( \hasstate(\sigma') \wand \wpre{\beta}{\Phi} \right)}
      {\wpre{\gVis_i(x, k)}{\Phi}}
    \end{mathpar}
    Use an accompanying ghost theory to recover local rules (later).
    \begin{mathpar}
      \inferrule[wp-read]{\heapCtx \\ \ell\pointsto \alpha \\ \later (\ell\pointsto \alpha \wand \wpre{\alpha}{\Phi})}{\wpre{(!\ell)}{\Phi}}
      \and \inferrule[wp-store]{\heapCtx \\ \ell\pointsto \alpha \\ \later (\ell\pointsto \beta \wand \Phi(\gRet()))}{\wpre{(\ell\gets \beta)}{\Phi}}
    \end{mathpar}
  \end{frame}

  \section{PRNG Effect Theory}

  \begin{frame}[allowframebreaks]{pseudo-randomness as an effect}
    \begin{block}{pseudo-randomness}
      Regarding a deterministic state machine as random/non-deterministic.
      \begin{itemize}
        \item Generator state \(S\)
        \item Output type \(T\)
        \item Next state \(S \to S\)
        \item Read out \(S \to T\)
        \item Seed \(S \to S \to S\)
      \end{itemize}
      For simplicity, take \(S = T = \mathbb{N}\)
    \end{block}
    \begin{block}{PRNG effect signature}
      \begin{itemize}
        \item Create a generator: \(\Ins_{\new}(X) = \mathbf{1}\) and \(\Outs_{\new}(X) = \iLoc\).
        \item Deletion a generator: \(\Ins_{\del}(X) = \iLoc\) and \(\Outs_{\del} = \mathbf{1}\).
        \item Add entropy: \(\Ins_{\seed}(X) = \iLoc \times \mathbb{N}\) and \(\Outs_{\seed}(X) = \mathbf{1}\).
        \item Next sample: \(\Ins_{\rand}(X) = \iLoc\) and \(\Outs_{\rand}(X) = \mathbb{N}\).
      \end{itemize}
    \end{block}
    \framebreak
    \begin{block}{PRNG effect reifier}
      \(\State(X) \eqdef \iLoc \fpfn \mathbb{N}\).
      \begin{align*}
        r_{\new}((), \sigma) & \eqdef \Some(\ell, \sigma[\ell\mapsto 0]) \text{ \(\ell \notin \dom(\sigma)\) is a fresh location}\\
        r_{\del}(\ell,\sigma) & \eqdef \begin{cases}
          \Some((), \mathrm{delete}(\sigma,\ell)) & \ell \in \dom(\sigma)\\
          \None & \ell \notin \dom(\sigma)\\
        \end{cases} \\
          r_{\seed}((\ell,s), \sigma) & \eqdef \begin{cases}
            \Some((), \sigma[\ell\mapsto \mix(n,s)]) & \ell \in \dom(\sigma) \land \sigma(\ell) = n\\
            \None & \ell \notin \dom(\sigma)\\
          \end{cases}\\
            r_{\rand}(\ell,\sigma) & \eqdef \begin{cases}
              \Some(\out(n), \sigma[\ell\mapsto \upd(n)]) & \ell \in \dom(\sigma) \land \sigma(\ell) = n\\
              \None & \ell \notin \dom(\sigma)\\
            \end{cases}
            \end{align*}
            Wrappers around \(\gVis\):
            \begin{itemize}
              \item \(\mathrm{PRNG\_NEW}(k : \iLoc \to \IT)\)
              \item \(\mathrm{PRNG\_DEL}(\ell : \iLoc)\)
              \item \(\mathrm{PRNG\_SEED}(\ell : \iLoc,n : \mathbb{N})\)
              \item \(\mathrm{PRNG\_GEN\_k}(\ell : \iLoc, k: \mathbb{N} \to \IT)\)
            \end{itemize}
          \end{block}
        \end{frame}

        \begin{frame}{we need a ghost theory}
          \(\hasstate(\sigma : \iLoc \fpfn \mathbb{N})\) is a monolithic.
          Devise a ghost theory to recover local rules.

          \begin{mathpar}
            \inferrule[wp-reify]
            {\hasstate(\sigma)\\
              r_i(x,\sigma) = \Some(\beta,\sigma')\\
            \later \left( \hasstate(\sigma') \wand \wpre{\beta}{\Phi} \right)}
            {\wpre{\gVis_i(x, k)}{\Phi}}
          \end{mathpar}

          \pause

          \begin{alertblock}{\(\mathrm{PRNG\_DEL}\) is evil}
            \begin{enumerate}
              \item Deleting a shared generator is not safe.
              \item Deletion is only safe under exclusive ownership.
              \item Controlled duplication --- fractional ownership
              \item \(\lambda_{\prng}\) substructural type system
            \end{enumerate}
            Dropped for simplicity
          \end{alertblock}
        \end{frame}

        \begin{frame}{Towards PRNG program logics}
          \begin{mathpar}
            \inferrule[gh-partial-dup]{ }{\knownprng(\ell)\vdash\knownprng(\ell) \star \knownprng(\ell)}
            \and \inferrule[wp-prng-new]
            {\prngCtx\\
            \later\later(\forall \ell.\; \knownprng(\ell) \wand \wpre{k\;\ell}{\Phi})}
            {\wpre{\mathrm{PRNG\_NEW}\;k}{\Phi}}
            \and
            \inferrule[wp-prng-gen]
            {\prngCtx\\
              \later \knownprng(\ell)\\
            \later\later \forall n.\; \Phi(\gRet(n))}
            {\wpre{\mathrm{PRNG\_GEN}\;\ell}{\Phi}}
            \and
            \inferrule[wp-prng-seed]
            {\prngCtx\\
              \later \knownprng(\ell)\\
            \later\later\Phi(\gRet(()))}
            {\wpre{\mathrm{PRNG\_SEED}\;\ell\;m}{\Phi}}
          \end{mathpar}
        \end{frame}

        \begin{frame}{PRNG ghost theory}
          \begin{mathpar}
            \inferrule[gh-store-alloc]{\sigma : \iLoc \fpfn \mathbb{N}}{\vdash \pvs \hasprngs(\sigma)}
            \and \inferrule[gh-loc-alloc]{\ell \notin \dom(\sigma) \\ n \in \mathbb{N}}{\hasprngs(\sigma) \vdash \pvs (\hasprngs(\sigma[\ell\easgn n])\star \knownprng(\ell))}
            \and \inferrule[gh-get]{ }{\hasprngs(\sigma) \star \knownprng(\ell)\vdash \hasprngs(\sigma) \star \exists n \in \mathbb{N}\;.\;\sigma(\ell) = n}
            \and \inferrule[gh-set]{n \in \mathbb{N}}{\hasprngs(\sigma) \star \knownprng(\ell)\vdash \pvs \hasprngs(\sigma[\ell\easgn n])}
          \end{mathpar}
          \pause
          Crafting a RA? Maybe two RAs:
          One authoritative RA is insufficient.
          \begin{mathpar}
            \inferrule
            {\forall \sigma'.\;
              \mathcal{V}(\sigma \cdot \sigma')
              \implies
            \mathcal{V}(\sigma[\ell\easgn n] \cdot \sigma')}
            {\AuthFull \sigma \Redu \AuthFull \sigma[\ell\easgn n]}
          \end{mathpar}
          What if \(\ell\) is already assigned in \(\sigma'\)?
          \vfill
          \pause
          \(\gamma_V\) for \(\textsf{exclR}(\textsf{gmapUR}(\iLoc,\mathbb{N}))\) and
          \(\gamma_K\) for \(\textsf{authR}(\textsf{gsetR}(\iLoc))\)
          \[
            \textsf{has\_prngs}(\sigma) \eqdef \ownGhost{\gamma_V}{\Excl(\sigma)} * \ownGhost{\gamma_K}{\AuthFull\dom(\sigma)}
            \qquad
            \knownprng(\ell) \eqdef \ownGhost{\gamma_K}{\AuthFrag(\{\ell\})}
          \]
        \end{frame}

        \begin{frame}{PRNG program \& proof}
          A naive rejection sampling program:
          \[
            \begin{array}{rcl}
              \syn{las\_vegas}(sd, init, test)
  & \eqdef &
  \begin{array}[t]{l}
    \mathrm{PRNG\_NEW}\;(\lambda r. \\
    \qquad \mathrm{ALLOC}\;(\gRet(init))\;(\lambda s. \\
    \qquad\quad \mathrm{PRNG\_SEED}\;r\;sd; \\
    \qquad\quad \gwhile\;(\mathrm{neg\_bool}\;(test \cbvap \mathrm{READ}\;s)) \\
    \qquad\qquad (\mathrm{LET}\;x = \mathrm{PRNG\_GEN}\;r\ \mathrm{in}\ \mathrm{WRITE}\;s\;x) \\
    \qquad\quad \mathrm{READ}\;s))
  \end{array}
            \end{array}
          \]
          A simple partial correctness spec:
          \begin{align*}
              & \persistently\left(
                \forall n.\; \wpre{\{test \cbvap \gRet(n)\}}
                {\beta.\; (\beta \equiv \gRet(1) \land f(n) = \mathbf{T}) \lor (\beta \equiv \gRet(0) \land f(n) = \mathbf{F})}
              \right) \\
            \star & \heapCtx \star \prngCtx \\
                  & \proves
                  \wpre{\syn{las\_vegas}(sd, init, test)}
                  {\beta.\; \exists n.\; \beta \equiv \gRet(n) \land f(n) = \mathbf{T}}
          \end{align*}
          Proof: as if we are doing a proof for \texttt{HeapLang}.
        \end{frame}

  \section{Denotational Semantics for lambda PRNG}

  \begin{frame}[allowframebreaks]{\(\lambda_\prng\) to GItree}
    \((e : \iExpr\ \Gamma)\ (\senv : \Gamma \to \IT) \to (\sem{e}_\senv : \IT)\)
    \begin{align*}
      \sem{\vunit}_\senv &\eqdef \gRet(())
      \quad \sem{n}_\senv \eqdef \gRet(n)
      \quad \sem{\loc}_\senv \eqdef \gRet(\loc) \\
      \sem{x}_\senv &\eqdef \senv(x) \\
      \sem{e_1 \oplus e_2}_\senv &\eqdef \gnatop{\oplus}(\sem{e_1}_\senv)(\sem{e_2}_\senv) \\
      \sem{\syn{if}\;e_0\;\syn{then}\;e_1\;\syn{else}\;e_2}_\senv & \eqdef \gif(\sem{e_0}_\senv)(\sem{e_1}_\senv)(\sem{e_2}_\senv) \\
      \sem{\eapp{e_1}{e_2}}_\senv &\eqdef \sem{e_1}_\senv \cbvap \sem{e_2}_\senv \\
      \sem{\recval{f}{x}{\expr}}_\senv
                                  &\eqdef \fix\;\lambda (\alpha : \latert \IT).\; 
                                  \gFun\left(\latert\!\circ\!\left(\lambda \alpha'\;v.\;
                                  \sem{\expr}_{\senv[x\mapsto v,\;f\mapsto \alpha']}\right)(\alpha)\right) \\
      \sem{\kw{NewPrng}}_\senv &\eqdef \mathrm{PRNG\_NEW}\;\gRet \\
      \sem{\kw{DelPrng}\;e}_\senv &\eqdef \getret(\mathrm{PRNG\_DEL})(\sem{e}_\senv) \\
      \sem{\kw{Rand}\;e}_\senv &\eqdef \getret(\mathrm{PRNG\_GEN})(\sem{e}_\senv) \\
      \sem{\kw{Seed}\;e_1\;e_2}_\senv &\eqdef \mathrm{SeedGit}\;(\sem{e_1}_\senv)\;(\sem{e_2}_\senv)
    \end{align*}

    (object language) recursion as (meta theory) guarded recursion. \(\recval{f}{x}{e}\)
    \begin{enumerate}
      \item use \(\fix\), construct \(\IT\) given \(\alpha : \latert \IT\)
      \item use \(\gFun\), construct \(\latert (\IT \to \IT)\)
      \item use \texttt{laterO\_map}, construct \(\IT \to \IT\)
      \item take input \(v : \IT\), construct \(\IT\)
      \item interpret \(e\) with \(f = \alpha\) and \(x = v\)
    \end{enumerate}

    \framebreak
    \((K : \iECtx\ \Gamma)\ (\senv : \Gamma \to \IT) \to (\sem{K}_\senv : \IT \to \IT)\)
    \begin{align*}
      \sem{\ctxhole}_\senv(\alpha) &\eqdef \alpha\\
      \sem{K\;e}_\senv(\alpha) &\eqdef \sem{K}_\senv(\alpha)\cbvap\sem{e}_\senv
      \quad \sem{v\;K}_\senv(\alpha) \eqdef \sem{v}_\senv\cbvap\sem{K}_\senv(\alpha)\\
      \sem{K \oplus e}_\senv(\alpha) &\eqdef \gnatop{\oplus}(\sem{K}_\senv(\alpha))(\sem{e}_\senv) \\
      \quad \sem{v \oplus K}_\senv(\alpha) &\eqdef \gnatop{\oplus}(\sem{e}_\senv)(\sem{K}_\senv(\alpha))\\
      \sem{\syn{if}\;K\;\syn{then}\;e_1\;\syn{else}\;e_2}(\alpha) & \eqdef \gif(\sem{K}_\senv(\alpha))(\sem{e_1}_\senv)(\sem{e_2}_\senv)\\
      \sem{\kw{DelPrng}\;K}_\senv(\alpha) &\eqdef \getret(\mathrm{PRNG\_DEL})(\sem{K}_\senv(\alpha)) \\
      \sem{\kw{Rand}\;K}_\senv(\alpha) &\eqdef \getret(\mathrm{PRNG\_GEN})(\sem{K}_\senv(\alpha)) \\
      \sem{\kw{Seed}\;K\;e}_\senv(\alpha) &\eqdef \mathrm{SeedGit}\;(\sem{K}_\senv(\alpha))\;(\sem{e}_\senv)\\
      \sem{\kw{Seed}\;v\;K}_\senv(\alpha) &\eqdef \mathrm{SeedGit}\;(\sem{v}_\senv)\;(\sem{K}_\senv(\alpha))\\
    \end{align*}
    \(\mathrm{SeedGit}\;(\gRet\,\ell)\;(\gRet\,n) \equiv \mathrm{PRNG\_SEED(\ell,n)}\).\\
    \(\mathrm{SeedGit}\;(-)\;\beta\) and \(\mathrm{SeedGit}\;\alpha_v\;(-)\) are GItree homomorphisms.
    \vfill
    \(\sem{K}_\senv\) is always a GItree homomorphism.\\
    The translation respects the correct order of evaluation.
  \end{frame}

  \begin{frame}[allowframebreaks]{Computational Adequacy}
     Binary logical relation \(\alpha \curlyeqprec e\): if \(e\) runs to a value \(v\), then \(\sem{e}\) runs to \(\sem{v}\).
     \begin{align*}
       \mathcal{V} & : \iType \to \mathrm{Rel}(\IT,\iVal) \\
       \mathcal{V}_{\tunit}(\alpha_v,v) &\eqdef \alpha_v \equiv \gRet(()) \land v \equiv \vunit \\
       \mathcal{V}_{\tnat}(\alpha_v,v)  &\eqdef \exists n.\; \alpha_v \equiv \gRet(n) \land v \equiv n \\
       \mathcal{V}_{\syn{prng}}(\alpha_v,v) &\eqdef \exists \loc.\;
       \alpha_v \equiv \gRet(\loc) \land v \equiv \loc \star \knownprng(\loc) \\
       \mathcal{V}_{\type_1 \to \type_2}(\alpha_v,v) &\eqdef \exists f.\;
       \alpha_v \equiv \gFun(f) \land
       \Box\forall \beta_v,u.\; \mathcal{V}_{\type_1}(\beta_v,u) \wand
       \mathcal{E}_{\type_2}(\gFun(f)\cbvap \beta_v,\; \eapp{v}{u})\\
       \mathcal{E} & : \iType \to \mathrm{Rel}(\IT,\iExpr) \\[0.6ex]
       \mathcal{E}_{\type}(\alpha,e) &\eqdef \forall \sigma.\;
       \hasstate(\sigma) \star \hasprngs(\sigma) \wand
       \wpre{\alpha}{\alpha_v.\; \exists v,\sigma'.\;
         \begin{array}[t]{ll}
        &(e,\sigma) \redu^\ast (v,\sigma')\\
           \star & \hasstate(\sigma')\\
           \star & \hasprngs(\sigma')\\
           \star & \mathcal{V}_{\type}(\alpha_v,v)
       \end{array}} \\[4ex]
       \mathcal{C} &: (\Gamma : \mathrm{Ctx}) \to \mathrm{Rel}(\Gamma\to \IT, \Gamma \to \iVal) \\
       \mathcal{C}_{\Gamma}(\senv,\env) &\eqdef \forall x\in\Gamma.\; \mathcal{V}_{\Gamma(x)}(\senv(x),\env(x)) \\
       \typedSem{\Gamma}{e}{\type} & \eqdef
       \forall \senv,\env.\;
       \mathcal{C}_{\Gamma}(\env,\senv) \wand
       \mathcal{E}_{\type}(\sem{e}_\senv,\; e[\env])
     \end{align*}
   \end{frame}
   \begin{frame}{Proof for \texttt{compat\_Seed}}
     \[
       \typedSem{\Gamma}{e_l}{\syn{prng}} \wand \typedSem{\Gamma}{e_s}{\tnat} \wand
       { \typedSem{\Gamma}{\kw{Seed}\;e_l\;e_s}{\tunit} }
     \]
     \begin{enumerate}[<+->]
       \item Intro \(\senv : \Gamma \to \IT\), \(\env : \Gamma \to \iVal\), and \(\mathcal{C}_{\Gamma}(\env,\senv)\). Show \(\mathcal{E}_{\tunit}(\sem{\kw{Seed}\,e_l\,e_s}_\senv, \kw{Seed}\,e_l\,e_s)\).
       \item Intro \(\sigma_0\) and \(\hasstate(\sigma_0) \star \hasprngs(\sigma_0)\). Show \(\wpre{\sem{\kw{Seed}\,e_l\,e_s}_\senv}{\ldots}\).
       \item The interpretation \(\sem{\kw{Seed}\;e_l\;e_s}_\senv\equiv \mathrm{SeedGit}\;(\sem{e_l}_\senv)\;(\sem{e_s}_\senv)\).
       \item Specialize \(e_l\) to \(\mathcal{E}_{\syn{prng}}(\sem{e_l}_\senv, \env(e_l))\).
         Specialize on \(\sigma_0\) to get \(\wpre{\sem{e_l}_\senv}{\ldots}\).
       \item \(\mathrm{SeedGit}\;(-);\beta\) is a homomorphism. Run \((\sigma_0,\sem{e_l}_\senv)\) to \((\sigma_1,\gRet\,l)\) with \(\hasstate(\sigma_1)\star \hasprngs(\sigma_1)\) and \(\knownprng(l)\).
       \item Specialize \(e_s\) to \(\mathcal{E}_{\tnat}(\sem{e_s}_\senv, \env(e_s))\).
         Specialize on \(\sigma_1\) to get \(\wpre{\sem{e_s}_\senv}{\ldots}\).
       \item \(\mathrm{SeedGit}\;(\gRet\,l);(-)\) is a homomorphism. Run \((\sigma_1,\sem{e_s}_\senv)\) to \((\sigma_2,\gRet\,n)\) with \(\hasstate(\sigma_2)\star \hasprngs(\sigma_2)\).
       \item Now \(\mathrm{SeedGit}\;(\gRet\,l)\;(\gRet\,n) \equiv \mathrm{PRNG\_SEED}(l,n)\). Show\\
         \(\wpre{\{\mathrm{PRNG\_SEED}(l,n)\}}{\alpha_v . \alpha_v \equiv \gRet() \star (\sigma_2,\syn{Seed}\;l\;n) \redus (\langle\rangle,\sigma_3) \star \cdots}\)
         with \(\hasstate(\sigma_3)\star\hasprngs(\sigma_3)\)
       \item Run GItree with \texttt{WP-PRNG-SEED}. Run \(\lambda_\prng\) with head reduction rule for \(\syn{Seed}\).
     \end{enumerate}
   \end{frame}

  \section{Discussions}
  \begin{frame}{Future work}
    \begin{block}{Sum up and Recap}
      \begin{enumerate}
        \item Background: Recursion i Iris, SKLR i Iris, GItrees, etc.
        \item PRNG effect theory: ghost theory and program logic.
        \item GItree-based interpretation of \(\lambda_\prng\).
        \item Logical relations for computational adequacy, a repertoire.
      \end{enumerate}
    \end{block}

    \pause
    \vfill

    \small
    \begin{itemize}
      \item \demph{\(\mathrm{PRNG\_DEL}\): frac own (ghost theory) and substructural type system (\(\lambda_\prng\) language)}
      \item \demph{GItrees models untyped effectful lambda calculus.
        Generalizing guarded delay monads for typed effectful languages?}
      \item Reify into (higher-order) state\&maybe monad \(\xlongrightarrow{}\) reify into prob and/or non-deter monad.
      \item Probabilistic (true) randomness \(\xlongleftarrow{}\) Pseudo-randomness \(\xlongrightarrow{}\) Non-determinism\footnote{\(\forall\) / \(\exists\). Must/May}.\\
        Same signature; different reifiers, ghost theory, and program logics.
      \item Two extremes: full control \(\ell \pointsto \alpha\) v.s. no control \(\knownprng(\ell)\)\\
        Partial knowledge e.g., last observed state of an external automaton.
      \item GItree-based denotations might be fully abstract.
      \item \(\wpre{\sem{e}}{\Phi}\) and \(\wpre{e}{\overline{\Phi}}\). (cf. Program logics \`a la carte)
      \item Bisimulation on GItrees (cf. \texttt{eq-up-to-tau} in ITrees).
    \end{itemize}
  \end{frame}


\end{document}



